\chapter{Reproductive weighting and the selected-parent distribution}
\label{selection:chapter}

A population fit asks how well a function describes individuals under
their existing frequencies. Selection asks how those frequencies change
when weighted by expected reproductive output. This chapter performs that
second calculation exactly for a Gaussian population and a log-quadratic
conditional mean. It then derives a separate error bound for ignoring
curvature. No inheritance rule is used: the distribution obtained here
describes parents weighted by their reproductive contributions.

\section{The exact model and its integrability condition}
\label{selection:model}

Let $\vect X=\vect A-\vect\mu\sim\mathcal N(\vect0,\mat G)$ with
$\mat G\succ0$. The singular case will be treated on the occupied
support. Assume in this chapter that conditional expected output is
exactly
\begin{equation}
W_A(\vect\mu+\vect x)
 =\exp\!\left(\ell_0+\vect b^{\mathsf T}\vect x
                   -\frac12\vect x^{\mathsf T}\bGamma\vect x\right),
\qquad \bGamma=-\mat H=\bGamma^{\mathsf T}.
\label{selection:log_quadratic_weight}
\end{equation}
This assumption is stronger than having a quadratic Taylor approximation
near the mean. It allows the selection integrals to use this function
over the whole occupied support. The matrix $\bGamma$ need not initially
be positive semidefinite. Positive semidefiniteness will be imposed when
deriving covariance contraction under stabilising curvature.

As derived in Proposition~\ref{conditioning:selection_weight}, the
reproductively weighted density is
\begin{equation}
p_{\mathrm{sel}}(\vect x)
 =\frac{p(\vect x)W_A(\vect\mu+\vect x)}{\overline W},
\qquad
\overline W=\int_{\R^n}p(\vect x)W_A(\vect\mu+\vect x)\,d\vect x.
\label{selection:density_definition}
\end{equation}
A probability density requires $0<\overline W<\infty$. Positivity of
the fitness function does not by itself ensure finiteness.

Multiplying the normal density by
Equation~\eqref{selection:log_quadratic_weight} gives
\begin{equation}
p(\vect x)W_A(\vect\mu+\vect x)
 =\frac{e^{\ell_0}}{(2\pi)^{n/2}\det(\mat G)^{1/2}}
    \exp\!\left(-\frac12\vect x^{\mathsf T}\mat J\vect x
                     +\vect b^{\mathsf T}\vect x\right),
\qquad \mat J=\mat G^{-1}+\bGamma.
\label{selection:combined_precision}
\end{equation}
The matrix $\mat J$ is the combined precision: it collects the quadratic
coefficient from population frequency and the quadratic coefficient from
expected output. ``Precision'' here means inverse covariance in a normal
density. Its sign determines whether the product can be normalised.

\begin{proposition}[Integrability of the Gaussian quadratic weight]
\label{selection:integrability}
For $\mat G\succ0$, finite $\ell_0$ and $\vect b$, and symmetric
$\bGamma$, the normalising integral in
Equation~\eqref{selection:density_definition} is finite if and only if
$\mat J=\mat G^{-1}+\bGamma\succ0$.
\end{proposition}
\begin{proof}
Diagonalise the symmetric matrix as
$\mat J=\mat O\operatorname{diag}(\lambda_i)\mat O^{\mathsf T}$,
where $\mat O$ is orthogonal. Set $\vect y=\mat O^{\mathsf T}\vect x$
and $\vect d=\mat O^{\mathsf T}\vect b$. The Jacobian has absolute
determinant one, and the nonconstant exponential becomes
\[
\prod_{i=1}^n\exp(-\lambda_i y_i^2/2+d_i y_i).
\]
Because it is nonnegative, integration can be performed one coordinate
at a time. If every $\lambda_i>0$, completing each scalar square gives
a finite Gaussian integral. If some $\lambda_i<0$, that coordinate has
a positive quadratic exponent and its integral diverges. If
$\lambda_i=0$, the integral is $\int_{\R}e^{d_i y_i}\,dy_i$; it
diverges whether $d_i=0$ or not. Thus strict positivity in every occupied
direction is necessary as well as sufficient.
\end{proof}

In particular, a zero eigenvalue at the boundary cannot be rescued by
putting a zero gradient in that direction. The integrand would then be
constant along an unbounded direction, which still has infinite integral.

\paragraph{Biological interpretation.}
Population frequencies can make a locally increasing expected-output
function integrable even if the fitness curvature is not stabilising in
every direction. Conversely, an assumed output function that increases
too fast relative to the Gaussian frequency decline gives infinite
population mean output. Such a model cannot define the selected group by
the proposed normalisation.

\section{Complete the square, including the normalising constant}
\label{selection:complete_square}

Assume henceforth that $\mat J\succ0$. Define
\begin{equation}
\mat V=\mat J^{-1},
\qquad \vect m=\mat V\vect b.
\end{equation}
Since $\mat J\vect m=\vect b$, expansion gives
\begin{align}
(\vect x-\vect m)^{\mathsf T}\mat J(\vect x-\vect m)
 &=\vect x^{\mathsf T}\mat J\vect x
   -\vect x^{\mathsf T}\mat J\vect m
   -\vect m^{\mathsf T}\mat J\vect x
   +\vect m^{\mathsf T}\mat J\vect m\nonumber\\
 &=\vect x^{\mathsf T}\mat J\vect x
   -2\vect b^{\mathsf T}\vect x
   +\vect b^{\mathsf T}\mat V\vect b.
\end{align}
Therefore
\begin{equation}
-\frac12\vect x^{\mathsf T}\mat J\vect x+\vect b^{\mathsf T}\vect x
 =-\frac12(\vect x-\vect m)^{\mathsf T}\mat J(\vect x-\vect m)
   +\frac12\vect b^{\mathsf T}\mat V\vect b.
\label{selection:completed_square}
\end{equation}
The two cross terms combine because $\mat J$ is symmetric; no commutation
of $\mat G$ and $\bGamma$ has been assumed.

To integrate the remaining centred quadratic, diagonalise $\mat J$ as
in the preceding proof. Translation by $\vect m$ and orthogonal rotation
leave the integration element unchanged. The product of one-dimensional
Gaussian integrals is
\begin{align}
\int_{\R^n}
 e^{-(\vect x-\vect m)^{\mathsf T}\mat J(\vect x-\vect m)/2}
 \,d\vect x
 &=\prod_i\sqrt{\frac{2\pi}{\lambda_i}}\nonumber\\
 &=(2\pi)^{n/2}\det(\mat J)^{-1/2}.
\end{align}
Substituting this into Equation~\eqref{selection:combined_precision}
gives the exact population mean output:
\begin{equation}
\overline W
 =e^{\ell_0}\det(\mat G)^{-1/2}\det(\mat J)^{-1/2}
      \exp\!\left(\frac12\vect b^{\mathsf T}\mat J^{-1}\vect b\right).
\label{selection:normaliser}
\end{equation}
Dividing by this constant cancels every term independent of $\vect x$.
The selected density is
\begin{equation}
p_{\mathrm{sel}}(\vect x)
 =\frac{\det(\mat J)^{1/2}}{(2\pi)^{n/2}}
   \exp\!\left[-\frac12(\vect x-\vect m)^{\mathsf T}
                       \mat J(\vect x-\vect m)\right].
\label{selection:normalised_density}
\end{equation}
This establishes the distribution, not only its first two moments.

\begin{proposition}[Selected-parent Gaussian moments]
\label{selection:selected_moments}
Under the exact Gaussian/log-quadratic model and positive combined
precision,
\begin{align}
\Gsel&=(\mat G^{-1}+\bGamma)^{-1},\nonumber\\
\vect\Delta_{\mathrm{full}}
 &=\E[\vect A\given\mathrm{sel}]-\vect\mu
   =\Gsel\vect b.
\label{selection:moment_formulas}
\end{align}
The selected breeding-value distribution is
$\mathcal N(\vect\mu+\vect\Delta_{\mathrm{full}},\Gsel)$.
\end{proposition}
\begin{proof}
Equation~\eqref{selection:normalised_density} is the normal density with
mean $\vect m$ and covariance $\mat J^{-1}$. Returning from centred
$\vect X$ to $\vect A=\vect\mu+\vect X$ adds $\vect\mu$ to the mean
and leaves the covariance unchanged.
\end{proof}

The covariance is independent of $\vect b$ within this exact family, while
the selected mean is not. This does not establish slope independence of
selected covariance for an arbitrary fitness surface or a non-Gaussian
population. Those cases need their own weighted-moment calculation.

\subsection{A second route through derivatives of the normaliser}
\label{selection:normaliser_derivatives}

With $\mat G$, $\bGamma$ and $\ell_0$ fixed, regard $\overline W$ as a
function of the coefficient vector $\vect b$. Differentiation under the
integral brings down a factor $X_i$, just as in the moment-generating
function calculation. Positive combined precision keeps these polynomially
weighted integrals finite. Hence
\begin{align}
\frac{\partial\log\overline W}{\partial b_i}
 &=\frac{\E[X_iW_A]}{\E[W_A]}
   =\E[X_i\given\mathrm{sel}],\nonumber\\
\frac{\partial^2\log\overline W}{\partial b_i\partial b_j}
 &=\E[X_iX_j\given\mathrm{sel}]
    -\E[X_i\given\mathrm{sel}]\E[X_j\given\mathrm{sel}].
\label{selection:derivative_moments}
\end{align}
These are derivatives with respect to a model coefficient, not spatial
derivatives with respect to breeding value. From
Equation~\eqref{selection:normaliser}, the only $\vect b$-dependent
term in $\log\overline W$ is
$\vect b^{\mathsf T}\mat J^{-1}\vect b/2$. Its gradient is
$\mat J^{-1}\vect b$ and its Hessian is $\mat J^{-1}$, recovering
both selected moments independently from the normalising expression.

\section{Stabilising curvature, saddle curvature and occupied support}
\label{selection:curvature_cases}

Define the symmetric covariance-weighted curvature
\begin{equation}
\mat T=\mat G^{1/2}\bGamma\mat G^{1/2}.
\label{selection:T}
\end{equation}
Multiplication shows
\begin{align}
\mat J
 &=\mat G^{-1/2}(\mat I+\mat T)\mat G^{-1/2},\nonumber\\
\Gsel
 &=\mat G^{1/2}(\mat I+\mat T)^{-1}\mat G^{1/2}.
\label{selection:precision_factorisation}
\end{align}
The inverse in the second line reverses the order of the factors in the
first. This factorisation works even when $\mat G$ and $\bGamma$ do
not commute. Congruence by an invertible matrix preserves positive
definiteness, so integrability is equivalent to
$\mat I+\mat T\succ0$. If $\tau_i$ are the eigenvalues of $\mat T$,
the condition is $\tau_i>-1$ for every occupied direction.

The normaliser can consequently also be written as
\begin{equation}
\overline W
 =e^{\ell_0}\det(\mat I+\mat T)^{-1/2}
     \exp\!\left(\frac12\vect b^{\mathsf T}\Gsel\vect b\right).
\label{selection:whitened_normaliser}
\end{equation}
To see this determinant identity explicitly,
$\det(\mat J)=\det(\mat G)^{-1}\det(\mat I+\mat T)$ in
Equation~\eqref{selection:precision_factorisation}; its covariance
determinant then cancels the one in
Equation~\eqref{selection:normaliser}.

\begin{corollary}[Covariance contraction under stabilising curvature]
\label{selection:covariance_contraction}
If $\bGamma\succeq0$, then
$\mat0\prec\Gsel\preceq\mat G$ for $\mat G\succ0$.
\end{corollary}
\begin{proof}
Positive semidefiniteness of $\bGamma$ implies $\mat T\succeq0$.
The eigenvalues of $(\mat I+\mat T)^{-1}$ are then
$1/(1+\tau_i)$, all positive and no greater than one. Therefore
\[
\mat0\prec(\mat I+\mat T)^{-1}\preceq\mat I.
\]
For any vector $\vect v$, premultiplication and postmultiplication by
$\mat G^{1/2}$ preserve the corresponding nonnegative quadratic forms.
Equation~\eqref{selection:precision_factorisation} gives the stated
covariance order.
\end{proof}
This order says that variance in every fixed linear combination of traits
does not increase. It does not require every individual covariance entry
to decrease numerically.

An indefinite $\bGamma$ can still give a valid selected Gaussian. For
example, take $\mat G=\mat I_2$ and
$\bGamma=\operatorname{diag}(0.4,-0.2)$. The combined precision is
$\operatorname{diag}(1.4,0.8)\succ0$, giving selected variances $5/7$
and $5/4$. One direction contracts and the other expands. With
$\vect b=(0.6,-0.4)^{\mathsf T}$, the mean shift is
$(3/7,-1/2)^{\mathsf T}$. Replacing the second curvature by $-1$ gives
a zero precision direction and a divergent normaliser, even if that
direction's gradient is also zero.

For singular $\mat G$, use the support construction in
Section~\ref{gaussian:singular_support}:
$\vect X=\mat U\vect Z$, with positive-definite covariance
$\mat\Lambda$ for $\vect Z$. Replace $\vect b$ by
$\mat U^{\mathsf T}\vect b$ and $\bGamma$ by
$\mat U^{\mathsf T}\bGamma\mat U$ in every integral and inverse.
The normalising density is then $r$-dimensional, where
$r=\operatorname{rank}(\mat G)$, and selected moments are mapped back
by $\mat U$. A gradient or Hessian component outside the occupied
support cannot affect this weighted distribution. Reweighting parents
also cannot create variation outside their existing support.

\paragraph{Biological interpretation.}
The selected covariance describes variation among contributing parents.
The stabilising-curvature result is a within-episode statement. It does
not specify covariance after recombination, segregation, mutation or
other processes of inheritance, and it is not a general claim that
population covariance must decline across generations.

\section{Ignoring curvature changes a different calculation}
\label{selection:linear_comparison}

If curvature is omitted, the reproductive weight is proportional to
$\exp(\vect b^{\mathsf T}\vect x)$, and the same square-completion
calculation with $\bGamma=\mat0$ gives
\begin{equation}
\vect\Delta_{\mathrm{lin}}=\mat G\vect b,
\qquad \mat G_{\mathrm{sel,lin}}=\mat G.
\label{selection:linear_shift}
\end{equation}
Using the population-fitted affine log function instead of the tangent
adds the intercept $\tr(\mat H\mat G)/2$. Its exponential is a common
positive factor and cancels from the selection normalisation. Thus both
choices give the same first-order parental mean shift here, despite
their different quality as fits on the log scale.

For one trait with initial variance $g>0$, curvature $H=-\gamma$, and
$1+\gamma g>0$, the exact formulas reduce to
\begin{equation}
g_{\mathrm{sel}}=\frac{g}{1+\gamma g},
\qquad
\Delta_{\mathrm{full}}=\frac{gb}{1+\gamma g},
\qquad \Delta_{\mathrm{lin}}=gb.
\label{selection:scalar_moments}
\end{equation}
If $b\ne0$, subtraction gives
\begin{equation}
\Delta_{\mathrm{lin}}-\Delta_{\mathrm{full}}
 =\gamma g\,\Delta_{\mathrm{full}},
\qquad
\frac{|\Delta_{\mathrm{lin}}-\Delta_{\mathrm{full}}|}
     {|\Delta_{\mathrm{full}}|}=|\gamma g|.
\label{selection:scalar_relative_error}
\end{equation}
For positive stabilising curvature, ignoring curvature overestimates the magnitude
of this scalar shift. For an integrable negative $\gamma$, it
underestimates it. In several traits the two shift vectors need not be
parallel, so a single coordinatewise percentage is not the general result.

\begin{example}[A share of $32/33$ and a displacement error of $1/2$]
\label{selection:97_50_example}
Take $b=2$, $H=-0.5$ and $g=1$. The variance components and their share
are
\[
\Vlin=4,\qquad \Vquad=\tfrac12(0.5)^2=0.125,\qquad
\Ag=\frac{4}{4.125}=\frac{32}{33}\simeq0.9697.
\]
The exact selected variance is $1/(1+0.5)=2/3$, and the exact shift is
$2/(1+0.5)=4/3$. Ignoring curvature gives a shift of $2$, so its error
relative to the exact shift is
\[
\frac{2-4/3}{4/3}=\frac{2/3}{4/3}=\frac12.
\]
The linear fit describes about $97\%$ of variation in the quadratic log
target, yet the resulting reproductive weighting overestimates the
selected-parent displacement by $50\%$.
\end{example}

For the same $g=1$ and $\gamma=0.5$, raising $b$ from $2$ to $10$
raises $\Vlin$ from $4$ to $100$ while leaving $\Vquad=0.125$.
The share becomes $100/100.125\simeq0.9988$. The shifts are now $10$
and $10/1.5$, but their relative discrepancy remains $1/2$. Increasing
the slope's magnitude is therefore not an accuracy guarantee.

Nor does high population $R^2$ guarantee agreement in every ranking.
In the first example, omit the irrelevant constant and write
$q(x)=2x-x^2/4$. At two upper-tail breeding values,
\[
q(4)=4>q(5)=3.75,
\]
whereas the fitted line $2x-1/4$ orders $5$ above $4$.
These are tail values of the specified standard normal population; the
example does not estimate the frequency of ranking reversals. It shows
that a variance-explained criterion is not an all-pairs ordering condition.

\paragraph{Biological interpretation.}
The frequency-weighted residual variance assigns little aggregate weight
to a rare part of the population. Reproductive weighting additionally
depends on its expected output, so the same part can contribute
differently to the selected group. The displayed example computes both
operations exactly. It makes no claim about offspring breeding values:
their distribution still requires a model of what parents transmit.

\section{Deriving the displacement-error bound}
\label{selection:error_bound_section}

We now compare the two shifts in a norm that accounts for the population's
covariance. Define
\begin{equation}
\norm{\vect v}_{\mat G^{-1}}
 =\sqrt{\vect v^{\mathsf T}\mat G^{-1}\vect v}
 =\norm{\mat G^{-1/2}\vect v},
\qquad \vect c=\mat G^{1/2}\vect b.
\label{selection:covariance_norm}
\end{equation}
This measures displacement after removing the covariance scale. A
displacement along a low-variance direction is large in this norm
relative to one of the same trait-coordinate length along a
high-variance direction. Unlike ordinary Euclidean length in arbitrarily
chosen trait units, the quadratic form is invariant when displacement
and covariance undergo the same invertible coordinate transformation:
\[
(\mat D\vect v)^{\mathsf T}
 (\mat D\mat G\mat D^{\mathsf T})^{-1}
 (\mat D\vect v)=\vect v^{\mathsf T}\mat G^{-1}\vect v.
\]

Let $\mat M=(\mat I+\mat T)^{-1}$, with $\mat T$ defined in
Equation~\eqref{selection:T}. The factorisation already proved gives
\begin{align}
\vect\Delta_{\mathrm{full}}
 &=\mat G^{1/2}\mat M\mat G^{1/2}\vect b
   =\mat G^{1/2}\mat M\vect c,\nonumber\\
\vect\Delta_{\mathrm{lin}}&=\mat G^{1/2}\vect c.
\end{align}
Multiplying by $\mat G^{-1/2}$ therefore gives
\begin{align}
\mat G^{-1/2}\vect\Delta_{\mathrm{full}}
 &=\mat M\vect c,\nonumber\\
\mat G^{-1/2}
 (\vect\Delta_{\mathrm{lin}}-\vect\Delta_{\mathrm{full}})
 &=(\mat I-\mat M)\vect c.
\end{align}
From $(\mat I+\mat T)\mat M=\mat I$, subtract $\mat M$ to obtain
$\mat I-\mat M=\mat T\mat M$. Thus the two whitened vectors are
\begin{equation}
\mat G^{-1/2}\vect\Delta_{\mathrm{full}}=\mat M\vect c,
\qquad
\mat G^{-1/2}(\vect\Delta_{\mathrm{lin}}-\vect\Delta_{\mathrm{full}})
 =\mat T\mat M\vect c.
\label{selection:whitened_error_identity}
\end{equation}
This step uses the inverse identity, not a presumed commutation of the
original covariance and curvature matrices.

\begin{theorem}[Relative selected-parent displacement bound]
\label{selection:relative_bound}
Assume the exact Gaussian/log-quadratic model, positive combined
precision, and $\vect c=\mat G^{1/2}\vect b\ne\vect0$. Then
\begin{equation}
\frac{\norm{\vect\Delta_{\mathrm{lin}}-
                  \vect\Delta_{\mathrm{full}}}_{\mat G^{-1}}}
     {\norm{\vect\Delta_{\mathrm{full}}}_{\mat G^{-1}}}
 \leq\norm{\mat T}_2
 \leq\norm{\mat T}_{\mathrm F}
 =\sqrt{2\Vquad}.
\label{selection:full_bound}
\end{equation}
For singular covariance the theorem applies on its occupied support,
with a nonzero gradient component there. Stabilising curvature is
sufficient but not necessary: the bound also holds for an integrable
indefinite curvature matrix.
\end{theorem}
\begin{proof}
Since $\mat M$ is invertible and $\vect c\ne\vect0$, the vector
$\vect y=\mat M\vect c$ is nonzero. Equation~\eqref{selection:whitened_error_identity}
turns the left-hand side into
\[
\frac{\norm{\mat T\vect y}}{\norm{\vect y}}.
\]
The spectral norm is defined as the supremum of this ratio over all
nonzero vectors, giving the first inequality.

For a symmetric matrix, diagonalise
$\mat T=\mat O\operatorname{diag}(\tau_i)\mat O^{\mathsf T}$.
Its spectral norm is $\max_i|\tau_i|$, while its squared Frobenius norm
is $\tr(\mat T^{\mathsf T}\mat T)=\sum_i\tau_i^2$. A largest squared
term cannot exceed the sum of all the squared terms, proving the second
inequality. Finally, trace cyclicity gives
\begin{align}
\norm{\mat T}_{\mathrm F}^2
 &=\tr(\mat T^2)\nonumber\\
 &=\tr(\mat G^{1/2}\bGamma\mat G\bGamma\mat G^{1/2})\nonumber\\
 &=\tr(\bGamma\mat G\bGamma\mat G)
   =\tr(\mat H\mat G\mat H\mat G)=2\Vquad,
\end{align}
where $\mat H=-\bGamma$ changes both curvature signs. Taking square
roots proves the final equality. Restricting every step to the
positive-variance support proves the singular version.
\end{proof}

If the gradient has zero component on the support, both shifts are zero.
Their absolute discrepancy is zero, but the displayed relative error is
$0/0$ and is undefined. It should not be assigned a finite relative
accuracy solely by continuity from a chosen gradient direction.

\subsection{Eigenbasis formula and equality conditions}
\label{selection:eigenbasis_error}

The bound can be read more closely by resolving the whitened gradient
into curvature eigenvectors. Let $c_i$ be its coefficients in the
orthonormal eigenbasis of $\mat T$. The exact selected displacement has
coefficients $y_i=c_i/(1+\tau_i)$, because integrability ensures
$1+\tau_i>0$. Therefore the squared relative error is exactly
\begin{equation}
\mathcal E^2
 =\frac{\displaystyle\sum_i
           \frac{\tau_i^2c_i^2}{(1+\tau_i)^2}}
         {\displaystyle\sum_i
           \frac{c_i^2}{(1+\tau_i)^2}}.
\label{selection:eigen_error}
\end{equation}
It is a weighted mean of the squared curvature eigenvalues, with weights
proportional to the squared components of the exact displacement.
If $\vect c$ lies along a single curvature eigenvector $k$, then
$\mathcal E=|\tau_k|$. Equality with the spectral-norm bound requires
the displacement, and equivalently the gradient, to lie entirely in
eigenspaces whose eigenvalues have maximum absolute value. Repeated
eigenvalues need not select a unique direction. Equality between the
spectral and Frobenius norms requires at most one nonzero eigenvalue.
In one trait both norm bounds are therefore exact, recovering
Equation~\eqref{selection:scalar_relative_error}.

This makes precise why fixed share need not mean fixed displacement
accuracy. If $\norm{\vect c}=1$ and
$\mat T=\operatorname{diag}(0.1,0.8)$, then every gradient orientation
has
\[
\Vlin=1,\qquad \Vquad=\tfrac12(0.1^2+0.8^2)=0.325,
\qquad \Ag=\frac1{1.325}.
\]
Gradients along the first and second eigenvectors nevertheless give
relative errors $0.1$ and $0.8$. The variance-share statement fixes a
gradient length and curvature spectrum. The selected-parent calculation
also depends on their relative orientation through the inverse factors
in Equation~\eqref{selection:eigen_error}.

\subsection{Why absolute curvature matters}
\label{selection:absolute_scale}

The upper bound involves $\sqrt{2\Vquad}$, not the fraction
$\Vquad/(\Vlin+\Vquad)$. A small fraction may result from a large
linear component rather than small absolute curvature. More strongly,
consider multiplying the nonconstant natural-log expected-output
function by $t>0$ while holding the population fixed. This replaces
$\vect b$ and $\mat H$ by $t\vect b$ and $t\mat H$. Both variance
components are multiplied by $t^2$, so their share is unchanged. In the
one-trait stabilising case, however, the relative displacement error
becomes $t\gamma g$. It can increase while the share stays exactly the
same.

This is a change in the expected-fitness relationship: up to a common
factor, expected output is raised to the power $t$. It is not a change
of trait units or merely a change of logarithm notation. Multiplying all
expected outputs by a common constant instead adds a log intercept,
leaving both the share and the normalised reproductive weights unchanged.

\paragraph{Biological interpretation.}
The share answers how much of a specified quadratic description is
represented linearly across the population. The bound answers how much
selected-parent displacement can change when curvature is omitted under
the exact distributional model. Reporting the two absolute components
alongside the share preserves the information needed for this separate
check. Extending the calculation to descendants still requires the
inheritance assumptions developed in the next chapter.

\section{Connection to executable checks}

The \codepath{selected_parent_moments} group in
\codepath{code/reproduce_result_checks.py} starts with
quadrature nodes from the original population and reweights them by
expected output. It compares the resulting moments with the formulas
above for stabilising and integrable saddle examples. The same file's
\codepath{orientation_and_absolute_scale} group checks equal-share
examples with different displacement errors. The scalar $32/33$ and
$50\%$ calculation is also integrated independently in
\codepath{tutorials/introduction/check_examples.py}.
These numerical checks assess implementations for specified cases; the
normalisation and bound for the whole stated model follow from the
derivations in this chapter.
